\scanpage{44}
des ``espaces stratifiés'' (en un sens qui
va être précisé plus bas) dont les
``pièces'' sont des $K(\pi,1)$, \deleted{cela devrait}\inserted{tels}
les espaces modulaires de Teichmüller
et leurs compactifiés, et les sous-espaces
\inserted{et voisinages tubulaires \unsure{parmi ceux-ci}}
qu'on peut leur associer grâce à leurs
stratifications canoniques, le type
d'homotopie doit pouvoir s'exprimer
en termes d'un système inductif
convenable (le plus souvent à ens. d'in-
dices fini) des groupoïdes fondamentaux
desdites pièces.

Il me semble d'ailleurs après-coup,
réflexion faite, que cette théorie ``d'intégration''
des types d'homotopie, que \unsure{nébuleusement}
je voyais comme liée au grand yoga des
relations entre structures homotopiques (relatives
sur un topos disons) et les $n$-groupoïdes,
est en fait relativement triviale, et
justiciable sans doute de techniques \unsure{bi-}
simpliciales des plus standart. Il me
faudra revenir là-dessus par la suite.
\mspage{43}
